% Generated draft: requires mathematical and source review.
\documentclass[11pt]{article}
\usepackage[margin=25mm]{geometry}
\usepackage{fontspec}
\setmainfont{Latin Modern Roman}
\usepackage{amsmath,amssymb,mathtools,unicode-math}
\setmathfont{Latin Modern Math}
\usepackage{enumitem,needspace,xurl,hyperref}
\setlength{\parindent}{0pt}
\setlength{\parskip}{5pt}
\setlength{\emergencystretch}{3em}
\newcommand{\N}{\mathbb{N}}
\newcommand{\Z}{\mathbb{Z}}
\newcommand{\Q}{\mathbb{Q}}
\newcommand{\R}{\mathbb{R}}
\newcommand{\C}{\mathbb{C}}
\newcommand{\F}{\mathbb{F}}
\newcommand{\A}{\mathbb{A}}
\newcommand{\PP}{\mathbb{P}}
\newcommand{\E}{\mathbb{E}}
\newcommand{\eps}{\varepsilon}
\newcommand{\dd}{\,\mathrm{d}}
\begin{document}
\section*{UnsolvedMath ID: 2306020}
\textbf{Source problem code:} catalogue entry
\subsection{Definitions and mathematical statement}
Let $\mathbb D=\{z\in\mathbb C:|z|<1\}$ and $\mathbb T=\partial\mathbb D$, and let $C\subset\mathbb D$ be the closed curve in the source. A univalent function is a holomorphic injective map $f:\mathbb D\to\mathbb C$. In the usual terminology, a convex closed curve means the boundary of a convex planar region; it does not mean that the one-dimensional curve itself is a convex subset.

\textbf{Problem.} Characterize the curves $C$ for which there exists a univalent $f$ on $\mathbb D$ such that both $f(C)$ and $f(\mathbb T)$ are convex curves.

The boundary expression requires clarification because $f$ is only defined on the open disk. If it is shorthand for convexity of the image domain, one precise interpretation asks for
\[
f(\mathbb D)\text{ convex},\qquad
f(C)=\partial K\text{ for a convex planar region }K.
\]
If instead $f(\mathbb T)$ means the actual pointwise image of the unit circle as a convex closed curve, an appropriate boundary extension of $f$ must be required and its continuity and injectivity specified. These are different formulations when no boundedness or boundary regularity is imposed. The source does not choose between them.

The source also says closed curve without explicitly requiring it to be a Jordan curve or to have any regularity. Such restrictions must be included in the requested necessary-and-sufficient conditions or recovered from context rather than silently assumed. The supplied summary reports the problem open as of the 2018 edition, with no definitive resolution located.
\subsection{Short English statement}
Determine when a conformal map of the disk can make an interior closed curve convex while also having a convex outer image. The source leaves the outer boundary interpretation and the class of curves unspecified.
\subsection{Sources}
\begin{itemize}
\item Hayman, \emph{Research Problems in Function Theory} (2018), Problem 6.20, \url{https://arxiv.org/abs/1809.07200}; source-TeX extraction, accessed 2026-07-29. The 2018 update was retained for status triage.
\item Local registry \texttt{data/unsolvedmath/problems.json}, permanent ID $2306020$, source code catalogue entry, \emph{Research Problems in Function Theory --- Problem 6.20}. The record labels this open and OPEN-TRIAGE, with NEEDS\_REVIEW source-status and rights qualifications. No further theorem citation or attribution is supplied.
\end{itemize}
\end{document}
